\section{Technische Verträge der prototypischen Realisierung}
\label{app:implementation_contracts}

Die prototypische Realisierung verwendet explizite technische Verträge, um
die im Architekturkonzept definierten Informations-, Provenienz- und
Autorisierungsgrenzen in überprüfbare Bedingungen der Implementierung zu
überführen. Die Verträge beschreiben insbesondere zulässige
Eingangszustände, erforderliche Identitäts- und Referenzbindungen,
Invarianten eines Verarbeitungsschritts sowie die Bedingungen für die
Weitergabe eines Ergebniszustands.

Tabelle~\ref{tab:implementation_contracts} fasst die für den
Transformationsprozess wesentlichen Verträge zusammen. Die Bezeichnungen
entsprechen den im Entwicklungsrepository verwendeten Contract-Namen.

\begin{longtable}{p{4.5cm}p{8.5cm}}
\caption{Wesentliche technische Verträge der prototypischen Realisierung}
\label{tab:implementation_contracts}\\
\hline
\textbf{Vertrag} & \textbf{Zentrale Festlegung} \\
\hline
\endfirsthead

\hline
\textbf{Vertrag} & \textbf{Zentrale Festlegung} \\
\hline
\endhead

Canonical Engineering Subject Contract &
Legt vor der Persona-basierten Interpretation eine eindeutige Menge
fachlicher Referenzgegenstände fest. Jeder Referenzgegenstand bleibt mit
seinen konkreten Vorkommen in der Engineering-Quelle verbunden. \\

Persona Subject Interpretation Contract &
Legt fest, dass jede konfigurierte Persona jeden festgelegten
Engineering Subject genau einmal interpretiert. Die Personas dürfen die
Menge der Referenzgegenstände dabei nicht eigenständig verändern. \\

Field-Level Subject Consensus Contract &
Definiert den regelbasierten Vergleich der Persona-Interpretationen für
strukturierte Merkmale und erhält abweichende Interpretationen sowie
Unsicherheiten für die nachfolgende menschliche Prüfung. \\

Subject-Centric Human Engineering Review Contract &
Legt fest, welche quellengebundenen Informationen, Persona-Ergebnisse und
Vergleichszustände der menschlichen fachlichen Prüfung für einen
Engineering Subject bereitgestellt werden. \\

Relationship Review Decision Persistence Contract &
Bindet menschliche Entscheidungen über vorgeschlagene fachliche Beziehungen
an die beteiligten Engineering Subjects und erhält frühere Entscheidungen
als nachvollziehbare Vorgängerzustände. \\

Finalization and Approved Engineering Information Contract &
Legt fest, unter welchen Bedingungen geprüfte Informationen als
autorisierte Engineering-Information für nachfolgende
Verarbeitungsschritte bereitgestellt werden dürfen. \\

Model Placement Review Contract &
Definiert die Entscheidungsgrenze zwischen fachlich autorisierter
Engineering-Information und ihrer strukturellen Einordnung in den
vorgegebenen Modellierungsrahmen. Die Entscheidung über die Platzierung ist
von der anschließenden Modellbildung getrennt. \\

\hline
\end{longtable}

Ergänzende Verträge regeln die Persistenz der dabei entstehenden
Zwischenzustände, ihre Einbindung in die menschliche Prüfoberfläche sowie
die Auswahl des jeweils gültigen Verarbeitungspfads. Diese Verträge dienen
vor allem der technischen Integration der beschriebenen
Architekturmechanismen und verändern deren fachliche Entscheidungsgrenzen
nicht.